\documentclass[11pt]{article}
\usepackage[T1]{fontenc}
\usepackage[utf8]{inputenc}
\usepackage{lmodern}
\usepackage[margin=1in]{geometry}
\usepackage{amsmath,amssymb,amsthm}
\usepackage{microtype}
\usepackage[colorlinks=true,linkcolor=blue,citecolor=blue,urlcolor=blue]{hyperref}
\newtheorem{theorem}{Theorem}[section]
\newtheorem{lemma}[theorem]{Lemma}
\newtheorem{proposition}[theorem]{Proposition}
\theoremstyle{remark}
\newtheorem{remark}[theorem]{Remark}
\newcommand{\R}{\mathbb R}
\newcommand{\Sph}{\mathbb S}
\newcommand{\eps}{\varepsilon}
\newcommand{\dist}{\operatorname{dist}}
\newcommand{\conv}{\operatorname{conv}}
\newcommand{\Span}{\operatorname{span}}
\newcommand{\A}{\mathcal A}
\title{Bounded flat transversals for fat convex bodies}
\author{Kaiwen Jin\thanks{\texttt{yui.ui.siki@gmail.com}}}
\date{}
\begin{document}
\maketitle

\begin{abstract}
Fix \(d\geq3\), \(1\leq k\leq d-2\), \(p\geq k+2\), and a fatness
parameter \(\rho\geq1\). We prove that a finite family of \(\rho\)-fat
compact convex bodies in \(\R^d\) can be intersected by a bounded number
of affine \(k\)-flats whenever every \(p\) members contain \(k+2\) members
with a common \(k\)-flat transversal. The bound depends only on
\(d,k,p,\rho\); the bodies may have arbitrary sizes and convex boundaries.
The missing rounding step is obtained by grouping bodies around members
of minimum radius, then reducing the dimension of their transversals
on the sphere. A finite family of perturbations turns the resulting
approximate transversals into exact ones at the scale of the smallest
body in each group.
\end{abstract}

\section{Introduction}
\label{sec:introduction}

Geometric transversal questions ask how a condition on small subfamilies
controls the number of probes needed to reach an entire family.
For affine subspaces of intermediate dimension, the shape of the sets
matters. Here we consider convex bodies with a fixed bound on their
eccentricity but with no restriction on their individual sizes.

A compact convex body \(K\subset\R^d\) is \(\rho\)-fat if there are a
center \(x_K\) and positive radii \(r_K,R_K\) such that
\begin{equation}
 B(x_K,r_K)\subseteq K\subseteq B(x_K,R_K),
 \qquad R_K\leq\rho r_K.
 \label{eq:fatness}
\end{equation}
All balls are closed. The radii in \eqref{eq:fatness} are allowed to vary
from one body to another. A \(k\)-flat is a \(k\)-dimensional affine
subspace, and a \(k\)-transversal of a subfamily is one \(k\)-flat meeting
every member of that subfamily.

Jung and P\'alv\"olgyi proved a fractional Helly theorem for
\(k\)-transversals of fat convex sets \cite[Theorem~5]{JP}.
Combined with the fractional transversal theorem of Alon, Kalai,
Matou\v{s}ek and Meshulam \cite{AKMM}, in the published formulation of
\cite[Theorem~22]{JP}, this bounds the fractional number of probes under
the local condition considered below.
Jung and P\'alv\"olgyi also proved the corresponding integer bound for
balls and for families having common inner and outer radius bounds
\cite[Theorem~7 and Corollary~8]{JP}.
Their Conjecture~1 asks for the extension to general fat convex sets.
The remaining question is whether a fractional cover can be replaced by
a bounded family of flats when the radii vary freely.

\begin{theorem}
\label{thm:main}
Let \(d\geq3\), \(1\leq k\leq d-2\), \(p\geq k+2\) be integers,
and let \(\rho\geq1\).
There is a constant \(C=C(d,k,p,\rho)\) with the following property.
If \(\mathcal F\) is a finite family of \(\rho\)-fat compact convex bodies
in \(\R^d\), and every \(p\) members of \(\mathcal F\) contain \(k+2\)
members with a common \(k\)-transversal, then at most \(C\) \(k\)-flats
together intersect every member of \(\mathcal F\).
\end{theorem}

Thus the theorem gives the finite-family form of
\cite[Conjecture~1]{JP}, with the same fatness convention and local
condition. The constant is independent of the number of bodies, their
locations, their global size ratio, and their boundary description
complexity.

The proof uses the fractional cover only at the end. We first prove a
uniform weak-net statement for a finite family of bodies that are each
hit with positive probability by a finite distribution of flats.
Bodies are grouped by their overlap with a body of minimum outer radius.
There are boundedly many groups, and within each group a fixed
positive amount of probe mass hits both the smallest body and any
specified member. Moving those probes through the center of the smallest
body allows a radial projection and decreases the dimension of the
transversals.

The movement introduces an absolute error equal to the radius of that
smallest body. Every other member of the group is at least as large.
Convexity and fatness then give a ball of comparable radius near each
approximate probe, so a fixed finite collection of perturbations reaches
all the original bodies. On the sphere this correction is made in
finitely many gnomonic charts. The argument retains tangential
incidences and allows new probes outside the support of the original
distribution. The resulting bounds are not optimized.

Section~\ref{sec:geometry} gives the grouping and geometric facts.
Section~\ref{sec:weaknets} proves the spherical and affine weak-net
statements and then Theorem~\ref{thm:main}.

\section{Grouping and geometric facts}
\label{sec:geometry}

All probability measures below have finite support, and expectations
are finite weighted sums. Families are indexed, so labels are retained
when bodies or projected images coincide.

\subsection{Grouping sets of positive measure}

The following elementary observation lets us use one smallest scale
within each group.

\begin{lemma}
\label{lem:grouping}
Let \(A_1,\ldots,A_n\) be measurable sets in a finite probability space
with measure \(\mu\), and suppose that \(\mu(A_i)\geq\eps>0\).
Assign a positive number \(R_i\) to each label.
There is a partition of the labels into at most
\[
 M(\eps)=\lceil2/\eps^2\rceil
\]
groups, with a distinguished label \(i_\ell\) in each group \(I_\ell\),
such that, for every \(i\in I_\ell\),
\[
 R_i\geq R_{i_\ell},
 \qquad \mu(A_i\cap A_{i_\ell})\geq\delta,
 \qquad \delta=\eps^2/2.
\]
\end{lemma}

\begin{proof}
Choose a label of minimum \(R_i\) among those remaining and remove as
its group every remaining label whose overlap with it is at least
\(\delta\). Repeat, breaking ties arbitrarily.
The minimum-radius condition is immediate, and every chosen label
removes itself.

If \(i_1,\ldots,i_t\) are the chosen labels, then their pairwise overlaps
are less than \(\delta\). For
\(Z=\sum_{\ell=1}^t1_{A_{i_\ell}}\), the nonnegativity of its variance
gives
\[
 t^2\eps^2\leq(\mathbb E_\mu Z)^2
 \leq\mathbb E_\mu Z^2
 \leq t+t(t-1)\delta.
\]
Hence \(t(\eps^2-\delta)\leq1-\delta\), which implies
\(t\leq2/\eps^2\).
\end{proof}

Conditioning \(\mu\) on \(A_{i_\ell}\) is legitimate because its measure
is positive. In that conditional probability, every group member has
overlap mass at least
\(\delta/\mu(A_{i_\ell})\geq\delta\).
The grouping itself always uses the original measure.

\subsection{Perturbing an affine transversal}

An approximate transversal at one fixed absolute scale can be replaced
by a finite family of exact transversals.

\begin{lemma}
\label{lem:euclidean-rounding}
Fix \(0\leq j<m\), \(\sigma\geq1\), \(a\geq0\), \(s>0\), and an affine
\(j\)-flat \(F\subset\R^m\).
Every family of \(\sigma\)-fat compact convex bodies with specified
outer radii \(R_K\geq s\) and
\(\dist(K,F)\leq as\) can be hit by at most
\begin{equation}
 Q(m-j,\sigma,a),\qquad
 Q(b,\sigma,a)=
 \bigl(2\lceil4\sigma\sqrt b(a+1)\rceil+1\bigr)^b
 \label{eq:Q}
\end{equation}
\(j\)-flats parallel to \(F\).
\end{lemma}

\begin{proof}
Translate the origin onto \(F\), and project orthogonally onto its normal
space of dimension \(b=m-j\). Each image \(A\) is convex and has the same
concentric inner and outer radius bounds as its original body.
Choose \(y\in A\) with \(|y|\leq as\), and let \(c\) be the projected
center. With
\(\lambda=s/(2R_K)\) and \(z=(1-\lambda)y+\lambda c\), convexity gives
\[
 B(z,\lambda r_K)\subseteq A,\qquad
 \lambda r_K\geq\frac{s}{2\sigma},\qquad
 |z-y|\leq\frac{s}{2}.
\]
Here \(\lambda\leq1/2\), and \(y\in B(c,R_K)\).
Thus \(|z|\leq(a+1/2)s\).

In an orthonormal coordinate system on the normal space, take grid
points with spacing \(h=s/(4\sigma\sqrt b)\) and coordinate indices
\(|q_\nu|\leq\lceil(a+1)s/h\rceil\).
A nearest grid point to \(z\) lies in this finite grid and has distance
at most \(\sqrt b\,h/2=s/(8\sigma)\) from \(z\).
It therefore belongs to \(A\).
There are at most \(Q(b,\sigma,a)\) grid points, and the corresponding
parallel flats hit every original body.
\end{proof}

The lower bound \(R_K\geq s\) is essential to this use of a common grid.
The error in the lemma is \(as\), not \(aR_K\).

\subsection{Spherical projections}

On the unit sphere \(\Sph^m\subset\R^{m+1}\), distances and cap radii
are geodesic. A proper spherical convex body is a compact set with
positive inradius, contained in an open hemisphere and closed under
the shorter great-circle arcs. We also allow the whole sphere.
Spherical \(\sigma\)-fatness means concentric cap bounds
\[
 B(x_K,r_K)\subseteq K\subseteq B(x_K,R_K),
 \qquad 0<r_K\leq R_K\leq\sigma r_K,\quad R_K\leq\pi.
\]
A great \(j\)-sphere is the intersection with a linear
\((j+1)\)-dimensional subspace. For \(j=0\) it is an antipodal pair.
We fix
\[
 \gamma=\pi/100.
\]

For \(w\in\Sph^m\), the gnomonic map to
\(P_w=\{z:w\cdot z=1\}\) is
\begin{equation}
 g_w(u)=\frac{u}{w\cdot u},\qquad
 g_w^{-1}(w+x)=\frac{w+x}{\sqrt{1+|x|^2}},
 \quad x\perp w.
 \label{eq:gnomonic}
\end{equation}
It is defined on the open hemisphere centered at \(w\) and maps
great-circle arcs to line segments.
Its inverse has differential norm at most \(1\).
On \(B(w,\pi/4)\), the differential of \(g_w\) has radial factor
\(\sec^2 t\) and transverse factor \(\sec t\), where \(t=\dist(u,w)\);
both are at most \(2\).
Since this cap is geodesically convex,
\begin{equation}
 \dist(g_w^{-1}(z),g_w^{-1}(z'))\leq|z-z'|,
 \qquad |g_w(u)-g_w(v)|\leq2\dist(u,v)
 \label{eq:gnomonic-bounds}
\end{equation}
in the respective domains.
Consequently a \(\sigma\)-fat body contained in that cap has a convex
image with concentric Euclidean radii \(r_K,2R_K\), and is
\(2\sigma\)-fat.

\begin{lemma}
\label{lem:spherical-fattening}
Suppose \(R_K<\gamma\) and \(0<s\leq R_K\), and let \(K^s\) be the
closed \(s\)-neighborhood of \(K\) on the sphere.
Then \(D=\conv_{\Sph}(K^s)\) is a compact \(\sigma\)-fat spherical body,
with concentric radii \(r_K+s,R_K+s\), and
\[
 D\subseteq K^{2s}.
\]
\end{lemma}

\begin{proof}
The neighborhood contains \(B(x_K,r_K+s)\) and is contained in
\(B(x_K,R_K+s)\). The latter cap has radius below \(2\gamma<\pi/4\)
and is convex, so it contains \(D\). Moreover
\[
 \frac{R_K+s}{r_K+s}\leq\frac{R_K}{r_K}\leq\sigma.
\]
Use \(g_{x_K}\) on this cap. Each point of \(g_{x_K}(K^s)\) is within
\(2s\) of the convex set \(g_{x_K}(K)\), by
\eqref{eq:gnomonic-bounds}. Taking convex hulls gives
\[
 g_{x_K}(D)\subseteq g_{x_K}(K)+B(0,2s).
\]
The inverse bound in \eqref{eq:gnomonic-bounds} proves
\(D\subseteq K^{2s}\). The Euclidean convex hull in this chart is compact,
which also proves compactness of \(D\).
\end{proof}

\begin{lemma}
\label{lem:pole-projection}
Let \(D\subset\Sph^m\) be a \(\sigma\)-fat spherical convex body with
specified outer radius \(R<\pi/4\), and let \(\omega\in\Sph^m\).
If \(D\) avoids both \(\omega\) and \(-\omega\), then
\[
 \Pi_\omega(y)=
 \frac{y-(y\cdot\omega)\omega}
 {|y-(y\cdot\omega)\omega|}
\]
maps \(D\) to a compact spherical convex
\(\pi^2\sigma\)-fat body in the equatorial sphere
\(\Sph^{m-1}\subset\omega^\perp\).
Great \(j\)-spheres through \(\omega\), for \(j\geq1\), project to
great \((j-1)\)-spheres, and their incidences with \(D\) lift exactly.
\end{lemma}

\begin{proof}
Write \(x,r,R\) for the concentric cap data, put
\(t=\dist(x,\omega)\), \(s_0=\min(t,\pi-t)\), and
\(w=\Pi_\omega(x)\).
Since the inner cap avoids both poles, \(r<s_0\).
For \(u<s_0\), the image of \(B(x,u)\) is the cap centered at \(w\)
of radius
\begin{equation}
 \alpha_u=\arcsin\!\left(\frac{\sin u}{\sin t}\right).
 \label{eq:cap-projection}
\end{equation}
To see this, fix a direction \(v\in\omega^\perp\) and consider
\(y=\sin\theta\,v+\cos\theta\,\omega\), \(0<\theta<\pi\).
Maximizing \(y\cdot x\) shows that this semicircle reaches \(B(x,u)\)
exactly when
\(\angle(v,w)\leq\arcsin(\sin u/\sin t)\).
The condition \(u<s_0\) excludes the two endpoint poles.

If \(R<s_0\), both projected cap bounds have the same center and
\[
 \frac{\alpha_R}{\alpha_r}
 \leq\frac{\pi}{2}\frac{\sin R}{\sin r}
 \leq\frac{\pi^2}{4}\frac{R}{r}
 \leq\frac{\pi^2}{4}\sigma.
\]
If \(R\geq s_0\), then
\[
 \alpha_r\geq\frac{\sin r}{\sin t}
 \geq\frac{2r}{\pi s_0}\geq\frac{2}{\pi\sigma}.
\]
Use the whole equatorial sphere as an outer cap of radius \(\pi\),
centered at \(w\). Its radius ratio is at most \(\pi^2\sigma/2\).
Both cases give the stated conservative bound.

The image is compact because \(D\) avoids the poles.
Projection preserves positive combinations, followed by normalization,
so the image is spherically convex.
For completeness, its Euclidean convex hull cannot contain zero:
a vanishing positive combination of projected directions, after dividing
the coefficients by the positive projection lengths, would lift to a
nonzero positive combination of points of \(D\) whose transverse
component vanishes. Normalizing that combination gives a pole in \(D\).
It cannot be zero because \(D\) lies in one cap of radius below
\(\pi/4\).
A minimum-norm point of the compact convex hull strictly separates it
from zero, so the image lies in an open hemisphere.

If a linear \((j+1)\)-space contains \(\omega\), its intersection with
\(\omega^\perp\) has dimension \(j\).
Conversely, spanning such a \(j\)-space with \(\omega\) restores the
original dimension. A point of \(D\) belongs to the lifted space
precisely when its projected direction belongs to the equatorial space.
\end{proof}

\subsection{Perturbing a great-sphere transversal}

\begin{lemma}
\label{lem:spherical-rounding}
Fix \(0\leq j<m\), \(\sigma\geq1\), \(s>0\), and a great
\(j\)-sphere \(F\subset\Sph^m\).
Every family of \(\sigma\)-fat spherical bodies with
\[
 s\leq R_K<\gamma,\qquad \dist(K,F)\leq2s
\]
can be intersected by at most \(B(m,j,\sigma)\) great \(j\)-spheres,
where \(B\) is independent of \(s\) and the family.
\end{lemma}

\begin{proof}
Fix a \(\gamma\)-net on \(F\) of cardinality \(A_j\), depending only
on \(j\) and the fixed \(\gamma\).
For a body \(K\), choose \(y\in K\), \(f\in F\) with
\(\dist(y,f)\leq2s\), and a net point \(w\) with
\(\dist(f,w)\leq\gamma\).
Every point of \(K\) lies within
\[
 2R_K+2s+\gamma<5\gamma<\pi/4
\]
of \(w\).

In the chart \(g_w\), \(F\) becomes an affine \(j\)-flat.
The image body is \(2\sigma\)-fat, has specified outer radius
\(2R_K\geq s\), and is within \(4s\) of that affine flat.
Lemma~\ref{lem:euclidean-rounding} gives at most
\(Q(m-j,2\sigma,4)\) parallel affine flats hitting all bodies
assigned to this chart.
Every affine \(j\)-flat in \(P_w\) spans a linear
\((j+1)\)-space with the origin, so its inverse image lies on a
great \(j\)-sphere. These lifted spheres hit the original bodies.
Thus we may take
\begin{equation}
 B(m,j,\sigma)=A_jQ(m-j,2\sigma,4).
 \label{eq:B}
\end{equation}
For \(j=0\), use the two points of \(F\) as its net; every chart point
lifts to an antipodal pair.
\end{proof}

\section{Weak nets and the bounded transversal theorem}
\label{sec:weaknets}

\subsection{Great-sphere transversals}

We first prove the spherical statement needed for the affine problem.
For a spherical body \(K\), write \(H_j(K)\) for the family of great
\(j\)-spheres meeting it.

\begin{theorem}
\label{thm:spherical-net}
For \(m\geq1\), \(0\leq j<m\), \(\sigma\geq1\), and
\(\eps\in(0,1]\), there is a finite
\(S(m,j,\sigma,\eps)\) such that the following holds.
If \(\mu\) is a finitely supported probability measure on great
\(j\)-spheres and \(\mathcal F\) is a finite family of
\(\sigma\)-fat spherical convex bodies satisfying
\(\mu(H_j(K))\geq\eps\) for all \(K\in\mathcal F\), then at most
\(S(m,j,\sigma,\eps)\) great \(j\)-spheres hit the whole family.
\end{theorem}

\begin{proof}
We induct on \(j\), for all \(m>j\).
Bodies with \(R_K\geq\gamma\) contain caps of radius at least
\(\gamma/\sigma\), so a fixed metric net of \(\Sph^m\) with mesh
\(\gamma/(2\sigma)\) hits all of them.
Extend its points to great \(j\)-spheres.
Let \(G(m,\sigma)\) be the cardinality of this fixed net, and remove
these large bodies.

Apply Lemma~\ref{lem:grouping} to the remaining incidence ranges and
their specified outer radii. There are at most \(M(\eps)\) groups.
In one group, let the distinguished body have center \(\omega\) and
outer radius \(s\). Every group member has \(s\leq R_K<\gamma\).
Put \(\delta=\eps^2/2\).

For \(j=0\), a positive overlap means that some antipodal pair hits
both the distinguished body and \(K\). Hence
\(\dist(K,\{\omega,-\omega\})\leq s\).
Lemma~\ref{lem:spherical-rounding} covers the group by at most
\(B(m,0,\sigma)\) antipodal pairs. This proves the base case.

Now let \(j\geq1\), and condition \(\mu\) on probes hitting the
distinguished body.
Each such great sphere \(F\) has distance at most \(s\) from \(\omega\).
Normalize the orthogonal projection of \(\omega\) onto its linear space,
obtaining \(v\in F\), and rotate in \(\Span(v,\omega)\) to send \(v\)
to \(\omega\), fixing the orthogonal complement.
The rotation angle is at most \(s\).
If \(v=\omega\), take the identity rotation.
Every point moves by at most \(s\): for a rotation of angle \(\theta\),
the dot product of a unit vector with its rotated image is at least
\(\cos\theta\).

For a group member put \(D_K=\conv_{\Sph}(K^s)\).
Lemma~\ref{lem:spherical-fattening} makes it \(\sigma\)-fat, with
outer radius \(R_K+s<2\gamma\), and gives \(D_K\subseteq K^{2s}\).
A probe hitting both \(K\) and the distinguished body becomes a great
sphere through \(\omega\) hitting \(D_K\).

Project the rotated probes onto \(\omega^\perp\).
This defines one finitely supported probability measure \(\nu\) on
great \((j-1)\)-spheres, independently of the choice of \(K\).
For each \(D_K\) avoiding the poles, Lemma~\ref{lem:pole-projection}
gives a \(\pi^2\sigma\)-fat projected body, and
\[
 \nu\bigl(H_{j-1}(\Pi_\omega D_K)\bigr)
 \geq\frac{\mu(H_j(K)\cap H_j(K_i))}{\mu(H_j(K_i))}
 \geq\delta,
\]
where \(K_i\) denotes the distinguished body.
The induction hypothesis supplies
\(S(m-1,j-1,\pi^2\sigma,\delta)\) equatorial probes.
Lift them through \(\omega\). They hit all the corresponding
\(D_K\), and also all \(D_K\) containing a pole.
If no projected bodies remain, choose one arbitrary great
\(j\)-sphere through \(\omega\).

Each \(K\) is now within \(2s\) of one lifted sphere.
Apply Lemma~\ref{lem:spherical-rounding} to the subfamily assigned
to each lifted sphere; all its members have \(R_K\geq s\).
This covers the group by boundedly many exact probes.

More explicitly, define the finite bounds recursively by
\begin{align}
 S(m,0,\sigma,\eps)
 &=G(m,\sigma)+M(\eps)B(m,0,\sigma),\notag\\
 S(m,j,\sigma,\eps)
 &=G(m,\sigma)+M(\eps)B(m,j,\sigma)
 \max\{1,S(m-1,j-1,\pi^2\sigma,\eps^2/2)\},
 \quad j\geq1.
 \label{eq:S}
\end{align}
The recursion terminates at \(j=0\).
\end{proof}

\subsection{Affine transversals}

Let \(\A_{d,k}\) be the space of affine \(k\)-flats in \(\R^d\), and
write \(H_k(K)=\{L\in\A_{d,k}:L\cap K\ne\varnothing\}\).

We use one radial projection fact.
If a compact convex \(\sigma\)-fat Euclidean body avoids the origin,
its radial image is a compact spherical convex
\(\pi\sigma\)-fat body.
Indeed, for its concentric data \(x,r,R\), put \(t=|x|>r\).
The inner ball projects to the cap centered at \(x/t\) with radius
\(\arcsin(r/t)\).
If \(t>R\), the corresponding outer radius is \(\arcsin(R/t)\),
and the ratio is at most \((\pi/2)R/r\).
If \(t\leq R\), the inner radius is at least \(r/t\geq1/\sigma\);
use an outer cap of radius \(\pi\) with the same center.
Radial projection preserves normalized positive combinations.
The image lies in an open hemisphere: a nearest point \(y\) of the
body to the origin satisfies \(y\cdot z\geq|y|^2>0\) for every
point \(z\) of the body.
Compactness follows from avoidance of the origin.

\begin{theorem}
\label{thm:affine-net}
Fix \(d\geq2\), \(1\leq k<d\), \(\rho\geq1\), and \(\eps\in(0,1]\).
There is a finite \(W(d,k,\rho,\eps)\) such that, for every finitely
supported probability measure \(\mu\) on \(\A_{d,k}\), every finite
family of \(\rho\)-fat compact convex bodies satisfying
\(\mu(H_k(K))\geq\eps\) is hit by at most \(W(d,k,\rho,\eps)\)
affine \(k\)-flats.
\end{theorem}

\begin{proof}
Apply Lemma~\ref{lem:grouping} to the incidence ranges with sizes \(R_K\).
For one group let the distinguished body have center \(a\) and
outer radius \(s\), and condition the measure on its hitting flats.
Translate each such flat through \(a\) by the shortest displacement,
whose length is at most \(s\).

Set \(D_K=K+B(0,s)\). It has concentric radii \(r_K+s,R_K+s\)
and remains \(\rho\)-fat.
Every original probe hitting both \(K\) and the distinguished body
becomes a flat through \(a\) hitting \(D_K\).
Take \(a\) as the origin and push the conditional measure to the
great \((k-1)\)-spheres given by these linear \(k\)-spaces.
This measure hits the radial image of each \(D_K\) avoiding \(a\)
with probability at least \(\delta=\eps^2/2\).
Those images are \(\pi\rho\)-fat by the radial fact just proved.

Theorem~\ref{thm:spherical-net} gives at most
\(S(d-1,k-1,\pi\rho,\delta)\) probes on \(\Sph^{d-1}\).
Their linear lifts through \(a\) also hit every \(D_K\) containing \(a\).
If every \(D_K\) contains \(a\), choose one arbitrary \(k\)-flat through it.
Each original \(K\) is at distance at most \(s\) from one of these
lifted flats, and all group members have \(R_K\geq s\).
Lemma~\ref{lem:euclidean-rounding}, with \(a=1\), replaces each lifted
flat by at most \(Q(d-k,\rho,1)\) exact probes.
Taking all groups, a valid bound is
\begin{equation}
 W(d,k,\rho,\eps)=M(\eps)Q(d-k,\rho,1)
 \max\{1,S(d-1,k-1,\pi\rho,\eps^2/2)\}.
 \label{eq:W}
\end{equation}
\end{proof}

\subsection{Proof of the main theorem}

The published fractional results give the following precise input.
For fixed \(d,k,p,\rho\) in Theorem~\ref{thm:main}, there is
\(T=T(d,k,p,\rho)\geq1\) such that every nonempty family under its
hypotheses has finitely many flats \(L_1,\ldots,L_N\) and nonnegative
weights \(w_1,\ldots,w_N\) satisfying
\begin{equation}
 \sum_{\nu=1}^N w_\nu\leq T,\qquad
 \sum_{\nu:L_\nu\cap K\ne\varnothing}w_\nu\geq1
 \quad(K\in\mathcal F).
 \label{eq:fractional-cover}
\end{equation}
To obtain this, use \cite[Theorem~5]{JP} on the flat-incidence ranges and
\cite[Theorem~22]{JP}, which credits the fractional transversal theorem
to Alon, Kalai, Matou\v{s}ek and Meshulam \cite{AKMM}.
In the latter theorem take \(q=k+2\). Its hypotheses are therefore the
original local condition and the published fat-body fractional Helly
property.
Finite support in \eqref{eq:fractional-cover} causes no restriction:
on a finite body family there are only finitely many realizable
incidence patterns. Choose one representative flat for each pattern
and collect its weights. The resulting finite cover problem attains
its optimum.

\begin{proof}[Proof of Theorem~\ref{thm:main}]
If \(|\mathcal F|<p\), one flat through a point of each body gives at
most \(p-1\) probes.
Otherwise normalize the weights in \eqref{eq:fractional-cover}
by their total. That total is positive, and every body has hitting
probability at least \(1/T\).
Apply Theorem~\ref{thm:affine-net} with \(\eps=1/T\).
Thus a valid choice is
\[
 C(d,k,p,\rho)=\max\{p-1,W(d,k,\rho,1/T(d,k,p,\rho))\}.
\]
The definitions \eqref{eq:Q}, \eqref{eq:B}, \eqref{eq:S},
and \eqref{eq:W} show that every bound is finite and depends only on
the displayed fixed parameters.
\end{proof}

\paragraph{AI-use statement.}
This manuscript was prepared with substantial assistance from OpenAI
Codex, including proof exploration, source checking, and drafting.
The mathematical claims are submitted for independent review.

\begin{thebibliography}{9}
\bibitem{AKMM}
N. Alon, G. Kalai, J. Matou\v{s}ek, and R. Meshulam,
\emph{Transversal numbers for hypergraphs arising in geometry},
Advances in Applied Mathematics \textbf{29} (2002), no.~1, 79--101.
\url{https://doi.org/10.1016/S0196-8858(02)00003-9}.
\href{https://web.math.princeton.edu/~nalon/PDFS/akmm6.pdf}{Author manuscript (21~May~2001).}

\bibitem{JP}
A. Jung and D. P\'alv\"olgyi,
\emph{k-Dimensional Transversals For Fat Convex Sets},
Discrete \& Computational Geometry (2026).
\url{https://doi.org/10.1007/s00454-026-00852-8}.
\end{thebibliography}
\end{document}
